%%%%%%%%%%%%%%%%%%%%%%%%%%%%%%%%%%%%%%%%%%%%%%%%%%%%%%%%%%%%%%%%%%%%%%
% Lớp 9 | Chương IV. Hệ thức lượng trong tam giác vuông | Tỉ số lượng giác của góc nhọn
% Dạng 9H4?11-1: Tính tỉ số lượng giác của góc nhọn
%%%%%%%%%%%%%%%%%%%%%%%%%%%%%%%%%%%%%%%%%%%%%%%%%%%%%%%%%%%%%%%%%%%%%%


% nguồn: dan_tu_AI (nhập 17/11/2026)
\begin{ex}%[9H4N11-1]%[ma:cdbab36a]
	Cho tam giác $ABC$ vuông tại $A$. Khẳng định nào sau đây đúng?
	\choice
	{$\sin B=\dfrac{AB}{BC}$}
	{\True $\sin B=\dfrac{AC}{BC}$}
	{$\sin B=\dfrac{AC}{AB}$}
	{$\sin B=\dfrac{BC}{AC}$}
	\loigiai{
		Với góc nhọn $B$: cạnh đối là $AC$, cạnh huyền là $BC$.\\
		Do đó $\sin B=\dfrac{\text{cạnh đối}}{\text{cạnh huyền}}=\dfrac{AC}{BC}$.
	}
\end{ex}

% nguồn: dan_tu_AI (nhập 17/11/2026)
% nguồn: dan_tu_AI (nhập 17/11/2026)
\begin{ex}%[9H4V11-1]%[ma:db225eed]
	Một chiếc thang dài $4\text{ m}$ dựa vào một bức tường thẳng đứng, thang tạo với mặt đất một góc $60^\circ$. Chân thang cách chân tường bao nhiêu mét?
	\choice
	{$2\sqrt3\text{ m}$}
	{\True $2\text{ m}$}
	{$8\text{ m}$}
	{$1\text{ m}$}
	\loigiai{
		Thang, tường và mặt đất tạo thành tam giác vuông với cạnh huyền là thang ($4\text{ m}$), góc giữa thang và mặt đất là $60^\circ$.\\
		Khoảng cách $d$ từ chân thang đến chân tường là cạnh kề góc $60^\circ$:
		\[ \cos 60^\circ=\dfrac{d}{4}\Rightarrow d=4\cdot\dfrac12=2\text{ (m)}. \]
	}
\end{ex}

% nguồn: dan_tu_AI (nhập 17/11/2026)
\begin{ex}%[9H4H11-1]%[ma:eb39fde5]
	Cho tam giác $ABC$ vuông tại $A$ có $AB=6\text{ cm}$, $AC=8\text{ cm}$. Các khẳng định sau đúng hay sai?
	\choiceTF
	{\True $BC=10\text{ cm}$}
	{\True $\sin B=\dfrac45$}
	{$\tan C=\dfrac43$}
	{$\sin B+\cos C=1$}
	\loigiai{
		\begin{itemchoice}
			\itemch Đúng. $BC=\sqrt{6^2+8^2}=10\text{ (cm)}$.
			\itemch Đúng. $\sin B=\dfrac{AC}{BC}=\dfrac{8}{10}=\dfrac45$.
			\itemch Sai. $\tan C=\dfrac{AB}{AC}=\dfrac68=\dfrac34$.
			\itemch Sai. $\cos C=\dfrac{AC}{BC}=\dfrac45$ nên $\sin B+\cos C=\dfrac45+\dfrac45=\dfrac85$.
		\end{itemchoice}
	}
\end{ex}

% nguồn: dan_tu_AI (nhập 17/11/2026)
% nguồn: dan_tu_AI (nhập 17/11/2026)
\begin{bt}%[9H4V11-1]%[ma:edb30078]
	Một máy bay cất cánh theo phương tạo với mặt đất một góc $15^\circ$. Sau khi bay được quãng đường $10\text{ km}$ theo phương đó, máy bay ở độ cao bao nhiêu kilômét so với mặt đất (làm tròn đến hàng phần trăm)?
	\shortans{2,59}
	\loigiai{
		Quãng đường bay là cạnh huyền ($10\text{ km}$), độ cao $h$ là cạnh đối của góc $15^\circ$.\\
		$\sin 15^\circ=\dfrac{h}{10}\Rightarrow h=10\sin 15^\circ\approx 2{,}59\text{ (km)}$.
	}
\end{bt}

% nguồn: dan_tu_AI (nhập 17/11/2026)
% nguồn: dan_tu_AI (nhập 17/11/2026)
\begin{bt}%[9H4C11-1]%[ma:29737687]
	Từ đỉnh một tòa nhà cao $20\text{ m}$, một người nhìn thấy một ô tô đang đỗ với góc hạ (góc giữa phương nằm ngang và tia nhìn xuống ô tô) là $30^\circ$. Ô tô chạy thẳng về phía chân tòa nhà rồi dừng lại; lúc này người đó nhìn thấy ô tô với góc hạ $60^\circ$. Hỏi ô tô đã chạy được bao nhiêu mét (làm tròn đến hàng phần mười)? Bỏ qua chiều cao của người và kích thước ô tô.
	\shortans{23,1}
	\loigiai{
		Gọi $A$ là đỉnh tòa nhà, $H$ là chân tòa nhà, $B$ và $C$ là vị trí ô tô lúc đầu và lúc sau.\\
		Góc hạ bằng góc tạo bởi tia nhìn với mặt đất (so le trong) nên $\widehat{ABH}=30^\circ$, $\widehat{ACH}=60^\circ$.\\
		Tam giác $ABH$ vuông tại $H$: $\tan 30^\circ=\dfrac{AH}{BH}\Rightarrow BH=\dfrac{20}{\tan 30^\circ}=20\sqrt3\text{ (m)}$.\\
		Tam giác $ACH$ vuông tại $H$: $\tan 60^\circ=\dfrac{AH}{CH}\Rightarrow CH=\dfrac{20}{\sqrt3}\text{ (m)}$.\\
		Quãng đường ô tô đã chạy:
		\[ BC=BH-CH=20\sqrt3-\dfrac{20}{\sqrt3}=\dfrac{40}{\sqrt3}\approx 23{,}1\text{ (m)}. \]
	}
\end{bt}
